% Generisano: uređuj beleske/sNNN.tex i naslove u manifestu.
\documentclass[11pt,a4paper]{article}
% Zajednički preambulum za dokument za čitanje; nije Beamer tema.
\usepackage[a4paper,margin=20mm,headheight=15pt]{geometry}
\usepackage{lmodern,fontspec}
\usepackage[serbian]{babel}
\setmainfont{Latin Modern Sans}
\setsansfont{Latin Modern Sans}
\setmonofont{Latin Modern Mono}
\usepackage{amsmath,amssymb,mathtools,graphicx,xcolor}
\usepackage{array,booktabs,tabularx,multirow,makecell,listings,fancyhdr}
\pagestyle{fancy}\fancyhf{}
\newcommand{\BeleskeTekuciID}{NEOZNACEN}
\fancyhead[L]{\small Beleške uz slajd \texttt{\expandafter\detokenize\expandafter{\BeleskeTekuciID}}}
\fancyfoot[R]{\small\thepage}
\setlength{\parindent}{0pt}\setlength{\parskip}{6pt}
\newwrite\BeleskeMapa
\immediate\openout\BeleskeMapa=\jobname.notes-map.tsv
\AddToHook{shipout/before}{%
  \immediate\write\BeleskeMapa{\BeleskeTekuciID\space\thepage}}
\AddToHook{enddocument/afterlastpage}{\immediate\closeout\BeleskeMapa}
\newcommand{\BeleskeID}[1]{%
  \def\BeleskeZadatiID{#1}%
  \ifx\BeleskeTekuciID\BeleskeZadatiID\else
    \PackageError{beleske}{Pogresan ID beleske: #1}{Proveri mapiranje slajda.}%
  \fi}
\newcommand{\BeleskaSlajda}[4]{%
  \clearpage\gdef\BeleskeTekuciID{#1}%
  {\Large\bfseries Slajd \textnormal{\texttt{\detokenize{#1}}}: #3\par}\medskip
  \includegraphics[page=#2,width=\linewidth]{\BeleskePrezentacija}\par\medskip
  \input{#4}}

\newcommand{\BeleskePrezentacija}{build/01_logicke_funkcije.pdf}
\begin{document}
\BeleskaSlajda{01-s001}{1}{Logičke funkcije}{beleske/s001.tex}
\BeleskaSlajda{01-s002}{2}{Digitalni sistem}{beleske/s002.tex}
\BeleskaSlajda{01-s003}{3}{Apstraktan i realan sistem}{beleske/s003.tex}
\BeleskaSlajda{01-s004}{4}{Alati za analizu i sintezu}{beleske/s004.tex}
\BeleskaSlajda{01-s005}{5}{Redna veza prekidača}{beleske/s005.tex}
\BeleskaSlajda{01-s006}{6}{Paralelna veza prekidača}{beleske/s006.tex}
\BeleskaSlajda{01-s007}{7}{Preklopnik i negacija}{beleske/s007.tex}
\BeleskaSlajda{01-s008}{8}{Dva preklopnika}{beleske/s008.tex}
\BeleskaSlajda{01-s009}{9}{Logički opis upravljanja liftom}{beleske/s009.tex}
\BeleskaSlajda{01-s010}{10}{Funkcija, operacija i prikazi}{beleske/s010.tex}
\BeleskaSlajda{01-s011}{11}{Binarna promenljiva}{beleske/s011.tex}
\BeleskaSlajda{01-s012}{12}{Sve unarne operacije}{beleske/s012.tex}
\BeleskaSlajda{01-s013}{13}{Negacija (NOT)}{beleske/s013.tex}
\BeleskaSlajda{01-s014}{14}{Bafer}{beleske/s014.tex}
\BeleskaSlajda{01-s015}{15}{Sve binarne operacije}{beleske/s015.tex}
\BeleskaSlajda{01-s016}{16}{Logičko I (AND)}{beleske/s016.tex}
\BeleskaSlajda{01-s017}{17}{Logičko ILI (OR)}{beleske/s017.tex}
\BeleskaSlajda{01-s018}{18}{Ekskluzivno ILI (EXOR)}{beleske/s018.tex}
\BeleskaSlajda{01-s019}{19}{Logičko NI (NAND)}{beleske/s019.tex}
\BeleskaSlajda{01-s020}{20}{Logičko NILI (NOR)}{beleske/s020.tex}
\BeleskaSlajda{01-s021}{21}{Ekskluzivno NILI (EXNOR)}{beleske/s021.tex}
\BeleskaSlajda{01-s022}{22}{Bulova algebra: aksiome}{beleske/s022.tex}
\BeleskaSlajda{01-s023}{23}{Bulova algebra: teoreme 1–5}{beleske/s023.tex}
\BeleskaSlajda{01-s024}{24}{Bulova algebra: teoreme 6–9}{beleske/s024.tex}
\BeleskaSlajda{01-s025}{25}{Troulazna I i ILI kola}{beleske/s025.tex}
\BeleskaSlajda{01-s026}{26}{Troulazno NI: pogrešna i ispravna kaskada}{beleske/s026.tex}
\BeleskaSlajda{01-s027}{27}{Četiri realizacije troulaznog NI kola}{beleske/s027.tex}
\BeleskaSlajda{01-s028}{28}{Troulazno EXILI i kaskada}{beleske/s028.tex}
\BeleskaSlajda{01-s029}{29}{Troulazno EXNILI i kaskada}{beleske/s029.tex}
\BeleskaSlajda{01-s030}{30}{Aktivan nivo: I i ILI}{beleske/s030.tex}
\BeleskaSlajda{01-s031}{31}{Aktivan nivo: NI i NILI}{beleske/s031.tex}
\BeleskaSlajda{01-s032}{32}{Aktivan nivo: EXILI i EXNILI}{beleske/s032.tex}
\BeleskaSlajda{01-s033}{33}{Aktivan nivo na ulazima}{beleske/s033.tex}
\BeleskaSlajda{01-s034}{34}{De Morganovi obrasci i dualna kola}{beleske/s034.tex}
\BeleskaSlajda{01-s035}{35}{Dualnost dvoulaznih i višeulaznih kola}{beleske/s035.tex}
\BeleskaSlajda{01-s036}{36}{Ispitivanje dualnosti EXILI i EXNILI}{beleske/s036.tex}
\BeleskaSlajda{01-s037}{37}{Primer realizacije samo NI kolima}{beleske/s037.tex}
\end{document}
